\documentclass{beamer}
\usepackage{moreverb} 
\usepackage{listings}
% imprimir
% \documentclass[handout]{beamer} 
% \usepackage{pgfpages}
% \pgfpagesuselayout{4 on 1}[a4paper,landscape,border shrink=5mm]

\mode<presentation> {
  \usetheme{Warsaw}
  \setbeamercovered{transparent}
}

\usebackgroundtemplate{\includegraphics[width=\paperwidth]{format/libresoft-bg.png}}
% \usepackage[spanish]{babel}
\usepackage[utf8]{inputenc}
\usepackage{graphics}
\usepackage{amssymb} % Simbolos matematicos

%\definecolor{libresoftgreen}{RGB}{162,190,43}
%\definecolor{libresoftblue}{RGB}{0,98,143}

%\setbeamercolor{titlelike}{bg=libresoftgreen}

%% Metadatos del PDF.
\hypersetup{  
  pdftitle={LaTeX in a nutshell},
  pdfauthor={Miguel Vidal},
  pdfcreator={GSyC/Libresoft},
  pdfproducer=PDFLaTeX,
  pdfsubject={Master on Free Software},
}
%%

\begin{document}

\title{LaTeX in a nutshell}
\subtitle{Master on Libre Software 2010}
\institute{mvidal@gsyc.urjc.es \\ GSyC/Libresoft -- URJC} 
\author{Miguel Vidal} 
%\date{\today}
\date{September 24, 2010}

\frame{
\maketitle
\begin{center}
\includegraphics[width=6cm]{format/gsyc-urjc}
\end{center}
}

\frame{
~
\vspace{4cm}

\begin{flushright}
{\footnotesize
(cc) 2010 Miguel Vidal. \\
  Some rights reserved. This work is licensed under a Creative Commons 
  Attribution-Share Alike 3.0 License, available at 
  \url{http://creativecommons.org/licenses/by-sa/3.0/}
 

\bigskip

}
\end{flushright}
}
%%

\begin{frame}
  \frametitle{Agenda}

%  \begin{itemize}[<+->]
  \begin{itemize}
    \item What is LaTeX
    \item References
  \end{itemize}

\end{frame}



%%%%%%%%%%%%%%%%%%%%%%%%%%%%%%%%%%%%%%%%%%%%%%%%%%%%%%%%%%%%%%%%%%%%%%%

\begin{frame}
\frametitle{How to pronounce and spell ``LaTeX''}

\begin{itemize}

\item ``\TeX'', ``\LaTeX'', or ``LaTeX'', no ``Latex''!
% \item Pronnounced \/lɑ:tɛx\/ in Spanish and \/lɑ:tɛx\/ or \/lei:tɛx\/ in English.
% Greek letters tau, epsilon, and chi
\item Pronounced /látej/ or /latek/ (English), no `latex'. 
\item It derives from the Ancient Greek: $\alert{\tau\epsilon\chi}\nu\eta$ (\textit{\alert{tej}né}: ``skill, art, technique'')
\item $\chi$: \textit{Ji} letter (in Spanish, fricativa velar sorda, ``j'' sound), \textit{Chi} (``kai'') in English.
\item IPA (International Phonetic Alphabet): [x] phonem 
\end{itemize}

\end{frame}

%%%%%%%%%%%%%%%%%%%%%%%%%%%%%%%%%%%%%%%%%%%%%%%%%%%%%%%%%%%%%%%%%%%%%%%

\begin{frame}
\frametitle{\TeX}

\begin{itemize}

\item \TeX{} was designed and mostly written by Donald Knuth (1977-1988)
\item \TeX{} understands about 300 low-level commands (``primitives''). Primitives are rarely used directly by users.
\item One of the most sophisticated digital typographical systems in the world.
\item Popular in academia, especially in mathematics, computer science, engineering, and physics. 
\item Functionality is provided by \alert{format files} (predumped memory images of \TeX{} after large macro collections have been loaded).

\end{itemize}

\end{frame}

%%%%%%%%%%%%%%%%%%%%%%%%%%%%%%%%%%%%%%%%%%%%%%%%%%%%%%%%%%%%%%%%%%%%%%%

\begin{frame}
\frametitle{\LaTeX}

\begin{itemize}

\item The most widely used \alert{format} is \LaTeX, originally developed by Leslie Lamport, which incorporates document styles for books, letters, slides, etc., 
\item \LaTeX{} was originally written in the early 1980s by Leslie Lamport. The current version is \LaTeX2e.
\item Set of macros from \TeX primitives.  
\item \LaTeX{} is free software distributed under the terms of the LaTeX Project Public License (LPPL), OSI-compliant.

\end{itemize}

\end{frame}

%%%%%%%%%%%%%%%%%%%%%%%%%%%%%%%%%%%%%%%%%%%%%%%%%%%%%%%%%%%%%%%%%%%%%%%

\begin{frame}
\frametitle{Why \LaTeX?}

\begin{itemize}

\item Typesetting, not ``word processing'': high quality final output
\item Accurate, precise output (device independent)
\item Separate content and styling
\item Modular, powerful and flexibility
\item Easy to make global changes
\item Encourage content reuse

\end{itemize}

\end{frame}

%%%%%%%%%%%%%%%%%%%%%%%%%%%%%%%%%%%%%%%%%%%%%%%%%%%%%%%%%%%%%%%%%%%%%%%

\begin{frame}
\frametitle{Caveats}

\begin{itemize}

\item Not WYSIWIG.
\item Hard learning curve.
\item Absolute space/positioning is tricky.
\item Need to be compiled.

\end{itemize}

\end{frame}




%%%%%%%%%%%%%%%%%%%%%%%%%%%%%%%%%%%%%%%%%%%%%%%%%%%%%%%%%%%%%%%%%%%%%%%

\begin{frame}[fragile] 
\frametitle{The simplest document}

\lstset{ %
language=TeX,
backgroundcolor=\color{green},
frame=[frame=trBL]
}

% \begin{boxedverbatim} 
\begin{lstlisting}[]
\documentclass{article}
\begin{document}
Hello World! {\textit This} is \LaTeX.
\end{document}
\end{lstlisting} 
% \end{boxedverbatim} 

\end{frame}

%%%%%%%%%%%%%%%%%%%%%%%%%%%%%%%%%%%%%%%%%%%%%%%%%%%%%%%%%%%%%%%%%%%%%%%

\begin{frame}[fragile]
\frametitle{Document workflow}

\alert{Editor} (\texttt{foobar.tex}) $\to$ \alert{\LaTeX{} processor} (\texttt{foobar.dvi}) $\to$ \alert{display} (viewer/screen) $\to$ \alert{drivers} (\texttt{foobar.ps}, printer)

\lstset{ %
language=TeX,
backgroundcolor=\color{green},
frame=[frame=trBL]
}

\begin{lstlisting}[]
_$ latex foobar.tex (`tex &latex foobar.tex')
_$ dvips -o foobar.ps foobar.dvi (ps output)
_$ pdflatex foobar.tex (pdf output)
_$ hevea foobar.tex  (html output)
\end{lstlisting}

\end{frame}


%%%%%%%%%%%%%%%%%%%%%%%%%%%%%%%%%%%%%%%%%%%%%%%%%%%%%%%%%%%%%%%%%%%%%%%

\begin{frame}
\frametitle{Document Structure}

Two main environments:

\begin{itemize}
\item \alert{Preamble:} commands and macros that affect the entire document.
\begin{itemize}
\item \alert{Top matter:} author, title, date, institution, etc. 
\end{itemize}
\item \alert{Document environment:} body text
\end{itemize}

\end{frame}

%%%%%%%%%%%%%%%%%%%%%%%%%%%%%%%%%%%%%%%%%%%%%%%%%%%%%%%%%%%%%%%%%%%%%%%

\begin{frame}[fragile]
\frametitle{Document Structure: Preamble}

\begin{itemize}
\item Everything from the start of the LaTeX source file until the \\begin\{document\} command
\item It normally contains global commands that affect the entire document.
\end{itemize}

\lstset{ %
language=TeX,
backgroundcolor=\color{green},
frame=[frame=trBL],
basicstyle=\footnotesize
}

\begin{lstlisting}[]
\documentclass[options]{class}
\usepackage[options]{package} 
\end{lstlisting}

\textit{class} (mandatory): book, article, report, letter, slides...
\textit{package} (optional): to utilize external macros

\end{frame}

%%%%%%%%%%%%%%%%%%%%%%%%%%%%%%%%%%%%%%%%%%%%%%%%%%%%%%%%%%%%%%%%%%%%%%%

\begin{frame}[fragile]
\frametitle{Top Matter}

\begin{itemize}
\item Title, date
\item Information about the authors, such as name, address, email etc.
\end{itemize}

\lstset{ %
language=TeX,
backgroundcolor=\color{yellow},
frame=[frame=trBL],
basicstyle=\footnotesize
}

\begin{lstlisting}[]
\documentclass[11pt,a4paper,oneside]{report}
\usepackage[utf8]{inputenc} % utf-8 encoding
\usepackage{amssymb}  % math symbols

\begin{document}
\title{How to Structure a LaTeX Document}
\author{Andrew Roberts}
\date{December 2004}
\maketitle
\end{document}\end{lstlisting}

\alert{class} (mandatory): book, article, report \\
\alert{package} (optional): to utilize external macros (\texttt{inputenc, babel, graphics, amssymb}...)

\end{frame}



%%%%%%%%%%%%%%%%%%%%%%%%%%%%%%%%%%%%%%%%%%%%%%%%%%%%%%%%%%%%%%%%%%%%%%%

\begin{frame}[fragile]
\frametitle{Document Structure: Body text}

\begin{itemize}
\item Abstract
\item Parts, chapters, sections, subsections, 
\item Appendices, Bibliography...
\end{itemize}

\lstset{ %
language=TeX,
backgroundcolor=\color{green},
frame=[frame=trBL]
}

\begin{lstlisting}[]
\begin{document}
  ... text mixed with local commands ...
\end{document}
\end{lstlisting}


\end{frame}



%%%%%%%%%%%%%%%%%%%%%%%%%%%%%%%%%%%%%%%%%%%%%%%%%%%%%%%%%%%%%%%%%%%%%%%


\begin{frame}[fragile] 
\frametitle{How to Structure a LaTeX Document}

LaTeX allows to structure documents with a variety of hierarchical constructs:

\lstset{ %
language=TeX,
backgroundcolor=\color{yellow},
frame=[frame=trBL],
basicstyle=\footnotesize
}


\begin{lstlisting}[]
\part{part title} 
\chapter{chapter title} 	% only books and reports
\section{section title} 	
\subsection{subsection title} 
\subsubsection{subsubsection title} 	
\end{lstlisting}


\end{frame}


%%%%%%%%%%%%%%%%%%%%%%%%%%%%%%%%%%%%%%%%%%%%%%%%%%%%%%%%%%%%%%%%%%%%%%%

\begin{frame}[fragile]
\frametitle{Font Styles}


\lstset{ %
language=TeX,
backgroundcolor=\color{yellow},
frame=[frame=trBL],
basicstyle=\footnotesize
}


\begin{lstlisting}[]
\textrm{...} % roman font
\textit{...} % italics
\textbf{...} % bold
\texttt{...} % monospace - teletype
\textsf{...} % sans serif 
\textsc{...} % small capitals
\end{lstlisting}

Example: 

\begin{lstlisting}[]
\LaTeX{} was \texttt{originally} written in \textbf{1984} 
by \textsc{Leslie Lamport} and has become the 
\textit{dominant} method for using \TeX.
\end{lstlisting}

\begin{block}{Example}
\LaTeX{} was \texttt{originally} written in \textbf{1984} by 
  \textsc{Leslie Lamport} and has become the 
  \textit{dominant} method for using \TeX.
\end{block}

\end{frame}

%%%%%%%%%%%%%%%%%%%%%%%%%%%%%%%%%%%%%%%%%%%%%%%%%%%%%%%%%%%%%%%%%%%%%%%


\begin{frame}[fragile]
\frametitle{Font Sizes}


\lstset{ %
language=TeX,
backgroundcolor=\color{yellow},
frame=[frame=trBL],
basicstyle=\footnotesize
}


\begin{lstlisting}[]
\tiny  
\scriptsize 
\footnotesize 
\small 	      
\normalsize 
\large 	
\Large 	
\LARGE 	
\huge 	
\Huge 	
\end{lstlisting}

(Related to font size in \texttt{\\documentclass})

Example: 

\begin{lstlisting}[]
\LaTeX{} was \texttt{originally} written in \textbf{1984} 
by \textsc{Leslie Lamport} and has become the 
\textit{dominant} method for using \TeX.
\end{lstlisting}

\begin{block}{Example}
\LaTeX{} was \texttt{originally} written in \textbf{1984} by 
  \textsc{Leslie Lamport} and has become the 
  \textit{dominant} method for using \TeX.
\end{block}

\end{frame}

%%%%%%%%%%%%%%%%%%%%%%%%%%%%%%%%%%%%%%%%%%%%%%%%%%%%%%%%%%%%%%%%%%%%%%%

\begin{frame}
\frametitle{Some special features}

\begin{itemize}

\item Text aligned
\item $n > 1$ blank lines and empty spaces: one line or one space
\item Start a new paragraph: \textbackslash\textbackslash
\item Hyphenate the word (exceptional cases): \texttt{man\textbackslash-u\textbackslash-script}
\item \texttt{\textbackslash newline}, \texttt{\textbackslash newpage}

\end{itemize}

\end{frame}



%%%%%%%%%%%%%%%%%%%%%%%%%%%%%%%%%%%%%%%%%%%%%%%%%%%%%%%%%%%%%%%%%%%%%%%

\begin{frame}[fragile]
\frametitle{Environments}


\lstset{ %
language=TeX,
backgroundcolor=\color{yellow},
frame=[frame=trBL],
basicstyle=\footnotesize
}

\begin{lstlisting}[]
\begin{environment name}
\end{environment name} 
\end{lstlisting}

Some environments: center, itemize, enumerate, figure, flushright, quotation, table...

Example:
\begin{lstlisting}[]
\begin{itemize}
\item Foo
\item Bar
\end{itemize}
\end{lstlisting}

\end{frame}




%%%%%%%%%%%%%%%%%%%%%%%%%%%%%%%%%%%%%%%%%%%%%%%%%%%%%%%%%%%%%%%%%%%%%%%

\begin{frame}[fragile] 
\frametitle{A simple document}

\lstset{ %
language=TeX,
backgroundcolor=\color{green},
frame=[frame=trBL],
basicstyle=\footnotesize
}


% \begin{boxedverbatim} 
\begin{lstlisting}[]
\usepackage[utf8]{inputenc}
\title{The beauty of \LaTeX}
\author{John Doe}
\date{\today}

\begin{document}
  \maketitle

% This is the comment body.
``Mathematical books and journals do not look as beautiful 
as they used to. It is not that their mathematical 
content is unsatisfactory, rather that the old and 
well-developed traditions of typesetting have become 
too expensive. Fortunately, it now appears that 
mathematics itself can be used to solve 
this problem.''\footnote{Donald E. Knuth, 1979.}

\end{document}

\end{lstlisting} 
% \end{boxedverbatim} 

\end{frame}


%%%%%%%%%%%%%%%%%%%%%%%%%%%%%%%%%%%%%%%%%%%%%%%%%%%%%%%%%%%%%%%%%%%%%%%


\begin{frame} 
\frametitle{References}

\begin{itemize}
\item \textsc{Lamport}, Leslie. \textit{LaTeX: A document preparation system: User's guide and reference}, Addison-Wesley, 1994.
\item \textsc{Sanguino-Botella}, Javier. \textit{Iniciación a \LaTeX2e: Un sistema para preparar documentos}, Addison-Wesley, 1997.
\item \textsc{VV.AA.} \textit{\LaTeX: Una imprenta en sus manos}, ADI, 2000.
\item CTAN: the authoritative collection of materials related to the TeX typesetting system. \url{http://www.ctan.org/}
\item Guide to the LaTeX markup language: \url{http://en.wikibooks.org/wiki/LaTeX/}
\item CervanTeX: Grupo de usuarios hispanohablantes de TeX: \url{http://www.cervantex.es/}
\end{itemize}

\end{frame}

%%%%%%%%%%%%%%%%%%%%%%%%%%%%%%%%%%%%%%%%%%%%%%%%%%%%%%%%%%%%%%%%%%%%%%%



\end{document}

Ejemplos tipografía vs word:
http://nitens.org/taraborelli/latex
